% Root and DE independent full proof audits passed on 2026-11-15.
\begin{statement}{C.39}{Proved}{Pentagon-trivial characters at five, seven and thirteen}
For each \(p\in\{5,7,13\}\) there is a continuous surjective
homomorphism
\[
              \epsilon_p:\widehat{\mathrm{GT}}_0
                                      \longrightarrow\mathbb F_p
\]
which vanishes on the full pentagon group and is nonzero
on global multiplier one.
In particular, if \(Q_{\rm GT}\) is the quotient of the full
gentle group by the closed normal subgroup generated by the
full pentagon group, then
\[
          Q_{\rm GT}^{\rm ab}\longrightarrow
             C_5\times C_7\times C_{13}
                       \times\prod_{j\geq1}C_2
                                                               \tag{C39.1}
\]
is a continuous surjection.
This gives characters of the gentle group that vanish on
the pentagon group, not nonzero characters of the pentagon
group itself.
\end{statement}
\begin{proof}
Use C.38 at top conjugator degree \(m=12\), and write
\(L=\bigoplus_{n=3}^{12}L^{(n)}\) for the rational grading
of the actual local multiplier-one image \(U\).
The exact calculation of C.33 shows
\[
       \dim S_{12}=7,\qquad
       \dim[S,S]_{12}=6.
\]
Its nonzero rational annihilator is generated by the
functional \(\ell=(-4,1,8,-4,-17,14,84)\) in the
retained basis.
Extend this functional to \(L\) by zero in every degree
other than twelve.
The argument of C.38 normalizes the actual compact image
of \(u\mapsto\ell(\log u)\) and gives an onto homomorphism
\[
       \Phi:U\longrightarrow\mathbb Z_p,\qquad
       \Phi(gug^{-1})=\lambda_p(g)^{12}\Phi(u).
                                                               \tag{C39.2}
\]
The scalar used to normalize this image is nonzero and
will not affect any vanishing assertion.

We must show vanishing on the actual pentagon image, not
only on homogeneous solutions of its first equation.
Let \(P\) be the image of the full pentagon group in
\(G=\mathcal G/\mathcal U_{13}\), and put \(P_1=P\cap U\).
It is a closed compact subgroup of the actual matrix group.
Define
\[
          L_P=\operatorname{span}_{\mathbb Q_p}
                                      \{\log h:h\in P_1\}.
\]
The finite BCH/Vandermonde proof in C.38 shows that \(L_P\)
is a Lie subalgebra.  The same explicit contraction proof,
using a basis of logarithms of elements of \(P_1\), shows
that \(P_1\) contains \(\exp\) of a sufficiently small full
lattice in \(L_P\).

Choose the actual Galois element \(g_0\) of local multiplier
\(1+p\) used in C.38.
It lies in \(P\), so \(L_P\) is stable under
\(A=\operatorname{Ad}(g_0)\).
The rational polynomial projector \(e_{12}\) for the
eigenvalue \((1+p)^{12}\) therefore preserves \(L_P\).
For any \(h\in P_1\), put
\[
                         X=e_{12}\log h .
\]
It belongs to \(L_P\cap L_{12}\).
For sufficiently large \(r\), the open exponential
neighborhood just proved gives
\(\exp(p^rX)\in P_1\).
This is an actual top-filtered element of the pentagon
image; all its lower conjugator terms vanish.
Its first pentagon equation is consequently the linear
weight-twelve pentagon equation, with homogeneous
conjugator \(p^rX\) under the identification with \(S_{12}\).

The exact rational row identity proved and executed in
C.35 gives
\[
 698922\,\ell
   =-192329R_1+472R_2-26842R_3-107926R_4-11870R_5.
                                                               \tag{C39.3}
\]
Every \(R_i\) vanishes on that actual leading pentagon
term.  Hence \(\ell(p^rX)=0\), and therefore
\(\ell(X)=0\) in \(\mathbb Q_p\).
This deduction is valid also when \(p\) divides the
integer in (C39.3): the identity is applied over
\(\mathbb Q_p\), before normalizing the actual compact
image or reducing it.
Since \(\ell\) is zero on all other eigenspaces,
\(\ell(\log h)=\ell(X)=0\).
Thus
\[
                              \Phi(P_1)=0.                     \tag{C39.4}
\]

For \(p=5,7,13\), the integer \(p-1\) divides twelve.
Reduction of (C39.2) modulo \(p\) is therefore invariant
under all of \(G\).
Extend it to a character of \(G\) using the vanishing
of \(H^2(\mathbb Z_p^\times,\mathbb F_p)\), and subtract
its cyclotomic Galois restriction as in C.38.
The resulting full gentle character is onto, since
the correction is zero on global multiplier one
and all local multiplier-one values have actual
complete full-group lifts.

It vanishes on \(P\).
Indeed for each \(g\in P\), choose a Galois element
\(\sigma\in P\) with the same local multiplier.
Then \(g\sigma^{-1}\in P_1\).
The corrected character is zero on \(\sigma\)
by construction and on \(P_1\) by (C39.4), so it
is zero on \(g\).
This proves the stated assertions for all three primes.
Their orders are pairwise coprime, and are also coprime
to the quadratic family of C.30.
Every finite product is consequently onto; compactness
gives (C39.1).
\end{proof}
